% =====================================================================
\section{Cơ sở lý thuyết}\label{sec:theory}
% =====================================================================

Phần này trình bày các khái niệm và định lý nền tảng cần dùng trong báo cáo.
Mọi định nghĩa được phát biểu kèm giải thích ý nghĩa, các định lý quan trọng
được chứng minh hoặc phác thảo chứng minh đầy đủ.

\subsection{Một số khái niệm chuẩn bị}\label{subsec:prelim}

\begin{definition}[Hệ trực chuẩn và ma trận trực giao]\label{def:orthonormal}
Cho không gian Euclide $\mathbb{R}^m$ với tích trong chuẩn
$\langle x, y\rangle = x^T y$. Hệ véc-tơ
$\{u_1, u_2, \dots, u_k\} \subset \mathbb{R}^m$ được gọi là \emph{trực chuẩn}
nếu
\[
    u_i^T u_j = \delta_{ij} =
    \begin{cases} 1 & \text{nếu } i = j,\\ 0 & \text{nếu } i \neq j.\end{cases}
\]
Ma trận $U \in \mathbb{R}^{m\times m}$ với các cột tạo thành hệ trực chuẩn
được gọi là \emph{ma trận trực giao}. Khi đó $U^T U = U U^T = I_m$ và
$U^{-1} = U^T$.
\end{definition}

\begin{definition}[Chuẩn Frobenius]\label{def:frobenius}
Với $A \in \mathbb{R}^{m\times n}$, chuẩn Frobenius được xác định bởi
\[
    \|A\|_F \;=\; \sqrt{\sum_{i=1}^{m}\sum_{j=1}^{n} a_{ij}^2}
            \;=\; \sqrt{\operatorname{tr}(A^T A)}
            \;=\; \sqrt{\operatorname{tr}(A A^T)}.
\]
Chuẩn này bất biến qua phép nhân với ma trận trực giao: nếu $Q$ trực giao thì
$\|Q A\|_F = \|A\|_F$ và $\|A Q\|_F = \|A\|_F$.
\end{definition}

\begin{proof}[Bất biến với phép trực giao]
Nếu $Q$ là ma trận trực giao thì
$\|Q A\|_F^2 = \operatorname{tr}((QA)^T (QA))
= \operatorname{tr}(A^T Q^T Q A) = \operatorname{tr}(A^T A) = \|A\|_F^2.$
Tương tự cho nhân bên phải.
\end{proof}

\begin{definition}[Hạng của ma trận]\label{def:rank}
Hạng $r = \operatorname{rank}(A)$ của ma trận $A \in \mathbb{R}^{m\times n}$
là số chiều của không gian cột (column space) của $A$, đồng thời cũng bằng số
chiều của không gian hàng. Luôn có $r \le \min(m, n)$.
\end{definition}

\subsection{Định lý phân tích giá trị kỳ dị}\label{subsec:svd-theorem}

\begin{theorem}[SVD đầy đủ]\label{thm:svd}
Với mọi ma trận $A \in \mathbb{R}^{m\times n}$ tồn tại các ma trận trực giao
$U \in \mathbb{R}^{m\times m}$, $V \in \mathbb{R}^{n\times n}$ và một ma trận
đường chéo suy rộng $\Sigma \in \mathbb{R}^{m\times n}$ sao cho
\begin{equation}\label{eq:svd}
    A \;=\; U\, \Sigma\, V^T,
\end{equation}
trong đó $\Sigma$ chỉ có các phần tử khác $0$ nằm trên đường chéo chính,
ký hiệu $\sigma_i = \Sigma_{ii}$ với $i = 1, 2, \dots, p$, $p = \min(m, n)$,
và thỏa mãn
\[
    \sigma_1 \;\ge\; \sigma_2 \;\ge\; \dots \;\ge\; \sigma_p \;\ge\; 0.
\]
Các số $\sigma_i$ được gọi là \emph{giá trị kỳ dị} của $A$; các cột của $U$
gọi là \emph{véc-tơ kỳ dị trái}, các cột của $V$ là \emph{véc-tơ kỳ dị phải}.
\end{theorem}

\begin{proof}[Phác thảo chứng minh dựa trên định lý phổ]
Ma trận $A^T A \in \mathbb{R}^{n\times n}$ là đối xứng và nửa xác định dương:
với mọi $x \in \mathbb{R}^n$, $x^T A^T A\, x = \|Ax\|_2^2 \ge 0$. Theo định lý
phổ cho ma trận đối xứng, tồn tại $V \in \mathbb{R}^{n\times n}$ trực giao và
$D = \operatorname{diag}(\lambda_1, \dots, \lambda_n)$ với
$\lambda_1 \ge \dots \ge \lambda_n \ge 0$ sao cho $A^T A = V D V^T$.

Đặt $\sigma_i = \sqrt{\lambda_i}$ với $i = 1, \dots, r = \operatorname{rank}(A)$.
Với $i \le r$, định nghĩa $u_i = (1/\sigma_i)\, A v_i$ (trong đó $v_i$ là cột
thứ $i$ của $V$). Kiểm tra trực tiếp:
\[
    u_i^T u_j = \frac{1}{\sigma_i \sigma_j} v_i^T A^T A v_j
              = \frac{\lambda_j}{\sigma_i\sigma_j} v_i^T v_j = \delta_{ij},
\]
nên $\{u_1, \dots, u_r\}$ trực chuẩn. Mở rộng thành cơ sở trực chuẩn của
$\mathbb{R}^m$ bằng quy trình Gram--Schmidt để được $U$ trực giao. Khi đó
$A V = U \Sigma$, hay tương đương $A = U\Sigma V^T$.
\end{proof}

\begin{remark}[Liên hệ với phổ]
Từ $A = U\Sigma V^T$ ta có
\[
    A^T A \;=\; V\, \Sigma^T \Sigma\, V^T, \qquad
    A A^T \;=\; U\, \Sigma \Sigma^T\, U^T.
\]
Vậy $\sigma_i^2$ chính là các trị riêng của $A^T A$ (và của $A A^T$),
$v_i$ là véc-tơ riêng của $A^T A$, $u_i$ là véc-tơ riêng của $A A^T$.
Đây là cơ sở để tính SVD bằng phương pháp lặp Lanczos/ARPACK hoặc phương
pháp Jacobi đối với $A^T A$.
\end{remark}

\subsubsection*{Ý nghĩa hình học của SVD}

Ánh xạ tuyến tính $x \mapsto Ax$ có thể được phân tích thành ba bước:
\begin{enumerate}
    \item Đổi cơ sở trực chuẩn qua $V^T$ (xoay/phản xạ trong $\mathbb{R}^n$);
    \item Co giãn theo từng trục bởi các hệ số $\sigma_i$;
    \item Đổi cơ sở trực chuẩn qua $U$ (xoay/phản xạ trong $\mathbb{R}^m$).
\end{enumerate}
Do đó, hình ảnh của hình cầu đơn vị qua $A$ là một \emph{ellipsoid} có các bán
trục bằng các giá trị kỳ dị $\sigma_i$, định hướng bởi các véc-tơ kỳ dị trái
$u_i$. \cref{fig:svd-geometry} minh họa ý tưởng này trong $\mathbb{R}^2$.

\begin{figure}[ht]
\centering
\begin{tikzpicture}[scale=1.0]
    % Left: unit circle
    \draw[->] (-2.2,0) -- (2.2,0) node[right]{$x_1$};
    \draw[->] (0,-2.2) -- (0,2.2) node[above]{$x_2$};
    \draw[thick, blue] (0,0) circle (1.5);
    \draw[->,thick,blue!70!black] (0,0) -- (1.5,0)
        node[below right]{$v_1$};
    \draw[->,thick,blue!70!black] (0,0) -- (0,1.5)
        node[above left]{$v_2$};
    \node at (0,-2.6) {Hình cầu đơn vị};

    % Middle arrow
    \node at (3.2,0) {$\xrightarrow{\; A \;}$};

    % Right: ellipse
    \begin{scope}[xshift=6.2cm]
    \draw[->] (-2.2,0) -- (2.6,0) node[right]{$y_1$};
    \draw[->] (0,-1.4) -- (0,1.4) node[above]{$y_2$};
    \draw[thick, red, rotate=20] (0,0) ellipse (2 and 0.8);
    \draw[->,thick,red!80!black,rotate=20] (0,0) -- (2,0)
        node[below right]{$\sigma_1 u_1$};
    \draw[->,thick,red!80!black,rotate=20] (0,0) -- (0,0.8)
        node[above left]{$\sigma_2 u_2$};
    \node at (0,-2.6) {Ảnh ellipsoid của $A$};
    \end{scope}
\end{tikzpicture}
\caption{Ý nghĩa hình học của SVD: ma trận $A$ biến đổi hình cầu đơn vị
trong $\mathbb{R}^n$ thành một ellipsoid trong $\mathbb{R}^m$. Các bán trục
chính có độ dài $\sigma_1 \ge \sigma_2 \ge \dots \ge \sigma_p$, định hướng
bởi các véc-tơ kỳ dị trái $u_i$.}\label{fig:svd-geometry}
\end{figure}

\subsection{Compact SVD và Reduced SVD}\label{subsec:compact-svd}

Gọi $r = \operatorname{rank}(A)$. Do
$\sigma_{r+1} = \sigma_{r+2} = \dots = \sigma_p = 0$, phân tích SVD đầy đủ có
thể được viết lại dưới dạng tổng dyadic
\begin{equation}\label{eq:dyadic}
    A \;=\; \sum_{i=1}^{p}\sigma_i u_i v_i^T
       \;=\; \sum_{i=1}^{r}\sigma_i u_i v_i^T.
\end{equation}

\begin{definition}[Compact SVD]\label{def:compact-svd}
Đặt
$U_r = [u_1 \mid u_2 \mid \dots \mid u_r] \in \mathbb{R}^{m\times r}$,
$\Sigma_r = \operatorname{diag}(\sigma_1, \dots, \sigma_r) \in \mathbb{R}^{r\times r}$,
$V_r = [v_1 \mid \dots \mid v_r] \in \mathbb{R}^{n\times r}$. Khi đó
\[
    A \;=\; U_r\, \Sigma_r\, V_r^T,
\]
được gọi là \emph{Compact SVD} của $A$. So với SVD đầy đủ, Compact SVD chỉ
giữ lại các thành phần thực sự đóng góp vào $A$, loại bỏ tất cả các cột
tương ứng với $\sigma_i = 0$.
\end{definition}

Khi $r \ll \min(m, n)$ (ma trận hạng thấp thật sự), Compact SVD đã tiết kiệm
đáng kể bộ nhớ: từ $\mathcal{O}(m^2 + n^2 + mn)$ phần tử của SVD đầy đủ xuống
$\mathcal{O}(r(m + n))$ phần tử.

\subsection{Định lý Eckart--Young: xấp xỉ hạng thấp tốt nhất}\label{subsec:eckart-young}

Định lý sau là kết quả nền tảng làm cơ sở cho mọi ứng dụng nén dữ liệu dựa
trên SVD.

\begin{theorem}[Eckart--Young 1936]\label{thm:eckart-young}
Cho ma trận $A \in \mathbb{R}^{m\times n}$ với SVD đầy đủ $A = U\Sigma V^T$
và $r = \operatorname{rank}(A)$. Với $1 \le k \le r$, đặt
\begin{equation}\label{eq:Ak}
    A_k \;=\; \sum_{i=1}^{k}\sigma_i u_i v_i^T \;=\; U_k\, \Sigma_k\, V_k^T,
\end{equation}
trong đó $U_k, \Sigma_k, V_k$ chứa $k$ cột/giá trị đầu tiên. Khi đó $A_k$ là
\emph{xấp xỉ hạng-$k$ tốt nhất} của $A$ theo chuẩn Frobenius, tức là
\begin{equation}\label{eq:eckart-young}
    \min_{B:\ \operatorname{rank}(B)\le k}\; \|A - B\|_F
       \;=\; \|A - A_k\|_F
       \;=\; \sqrt{\sum_{i=k+1}^{r}\sigma_i^2}.
\end{equation}
\end{theorem}

\begin{proof}[Phác thảo chứng minh]
Giả sử $B \in \mathbb{R}^{m\times n}$ có $\operatorname{rank}(B) \le k$. Khi
đó không gian zero (null space) của $B$ có chiều ít nhất $n - k$, nên giao với
không gian con $V_{k+1} = \operatorname{span}\{v_1, \dots, v_{k+1}\}$ (có chiều
$k+1$) phải có chiều ít nhất $1$. Lấy $w \in V_{k+1} \cap \ker(B)$ với
$\|w\|_2 = 1$, biểu diễn $w = \sum_{i=1}^{k+1}\alpha_i v_i$ với
$\sum \alpha_i^2 = 1$. Khi đó:
\[
    \|A - B\|_2^2 \;\ge\; \|(A-B) w\|_2^2 \;=\; \|A w\|_2^2
              \;=\; \sum_{i=1}^{k+1}\alpha_i^2 \sigma_i^2 \;\ge\; \sigma_{k+1}^2,
\]
chứng minh kết quả với chuẩn spectral. Đối với chuẩn Frobenius, lập luận sử
dụng tính bất biến trực giao và một bài toán tối ưu lồi trên không gian các
hệ số chéo (chi tiết xem Trefethen \& Bau \cite{trefethen1997}).
\end{proof}

\begin{remark}[Ý nghĩa thực tiễn]
\cref{thm:eckart-young} khẳng định: trong tất cả các ma trận hạng không vượt
quá $k$, $A_k$ làm tối thiểu sai số $\|A - B\|_F$, với cực tiểu chính bằng
căn bậc hai của tổng bình phương các giá trị kỳ dị bị bỏ qua. Đây là lý do mọi
phương pháp nén dữ liệu dựa trên SVD đều quy về việc \emph{giữ lại $k$ thành
phần kỳ dị lớn nhất}.
\end{remark}

\subsection{Năng lượng dữ liệu và tiêu chuẩn chọn $k$}\label{subsec:energy}

\begin{proposition}[Phân rã năng lượng]\label{prop:energy}
Với $A = U\Sigma V^T$, ta có
\begin{equation}\label{eq:energy-decomp}
    \|A\|_F^2 \;=\; \sum_{i=1}^{r}\sigma_i^2.
\end{equation}
\end{proposition}

\begin{proof}
Từ \cref{def:frobenius} và tính bất biến chuẩn Frobenius qua phép trực giao:
$\|A\|_F = \|U\Sigma V^T\|_F = \|\Sigma\|_F$. Vì $\Sigma$ chỉ có các phần tử
khác $0$ trên đường chéo nên
$\|\Sigma\|_F^2 = \sum_{i=1}^{p}\sigma_i^2 = \sum_{i=1}^{r}\sigma_i^2$
(do $\sigma_i = 0$ với $i > r$).
\end{proof}

\begin{definition}[Năng lượng dữ liệu]\label{def:energy}
Đại lượng $\|A\|_F^2$ được gọi là \emph{tổng năng lượng} (total energy) của
ma trận dữ liệu $A$. Mỗi $\sigma_i^2$ là \emph{năng lượng} đóng góp bởi thành
phần kỳ dị thứ $i$. Tỉ số
$\sigma_i^2 / \|A\|_F^2$ chính là tỉ trọng năng lượng của thành phần thứ $i$.
\end{definition}

Trên thực nghiệm với dữ liệu thực, dãy $(\sigma_i^2)$ thường suy giảm rất
nhanh theo $i$ -- đặc biệt với dữ liệu có cấu trúc latent (như ảnh tự nhiên,
ma trận tương quan). Đây chính là cơ sở để nén dữ liệu khả thi: một số ít
giá trị kỳ dị đầu tiên đã chứa phần lớn năng lượng.

\subsubsection*{Tiêu chuẩn chọn $k$ theo tỉ lệ năng lượng}

Cho ngưỡng $\eta \in (0, 1)$ (ví dụ $0.90$, $0.95$, $0.99$), ta chọn $k$ nhỏ
nhất thỏa mãn
\begin{equation}\label{eq:energy-criterion}
    \frac{\sum_{i=1}^{k}\sigma_i^2}{\sum_{i=1}^{r}\sigma_i^2} \;\ge\; \eta.
\end{equation}

\begin{corollary}[Cận sai số tương ứng]\label{cor:error-bound}
Với $k$ chọn theo \eqref{eq:energy-criterion}, sai số tương đối bình phương
được kiểm soát:
\[
    \frac{\|A - A_k\|_F^2}{\|A\|_F^2}
       \;=\; 1 - \frac{\sum_{i=1}^{k}\sigma_i^2}{\sum_{i=1}^{r}\sigma_i^2}
       \;\le\; 1 - \eta.
\]
\end{corollary}

\begin{proof}
Áp dụng \cref{thm:eckart-young} cho tử số và \cref{prop:energy} cho mẫu số:
\[
    \frac{\|A - A_k\|_F^2}{\|A\|_F^2}
       \;=\; \frac{\sum_{i>k}\sigma_i^2}{\sum_{i=1}^{r}\sigma_i^2}
       \;=\; 1 - \frac{\sum_{i \le k}\sigma_i^2}{\sum_{i=1}^{r}\sigma_i^2}
       \;\le\; 1 - \eta. \qedhere
\]
\end{proof}

\cref{cor:error-bound} cho biết: chọn $\eta = 0.95$ \emph{đảm bảo} sai số bình
phương tương đối không vượt quá $5\%$. Đây là một mệnh đề \emph{định lượng và
có thể kiểm chứng}, không phải một heuristic.

\subsection{Tính khả thi tính toán}\label{subsec:complexity-svd}

Mặc dù \cref{thm:eckart-young} cho ta xấp xỉ tốt nhất, việc \emph{tính} SVD
đầy đủ trên ma trận lớn rất tốn kém:

\begin{itemize}
    \item Thuật toán SVD chuẩn (Golub--Reinsch, LAPACK \texttt{GESVD}) có độ
    phức tạp $\mathcal{O}(mn\min(m,n))$ và yêu cầu lưu cả ma trận trung gian
    dưới dạng dày.
    \item Trên ma trận thưa, áp ngay GESVD đòi hỏi \emph{dày hóa} ma trận,
    phá vỡ cấu trúc thưa và làm tăng bộ nhớ lên cấp $mn$ phần tử.
    \item Với $m, n \sim 10^5$ (như MovieLens 25M), $mn \sim 10^{10}$ phần
    tử dày -- vượt xa giới hạn RAM tiêu chuẩn.
\end{itemize}

Đây là động lực cho Randomized SVD: tính nhanh \emph{xấp xỉ} cho $k$ giá trị
kỳ dị lớn nhất, với chi phí tỷ lệ thuận với $k$ và tận dụng được cấu trúc thưa
của ma trận đầu vào. \cref{sec:algorithm} sẽ trình bày chi tiết.
